% Compatibility definitions used by the thesis chapters migrated from the
% previous document class. Presentation is left to USCthesis.

\DeclareMathOperator*{\argmax}{arg\,max}
\DeclareMathOperator*{\argmin}{arg\,min}
\renewcommand{\binom}[2]{\left(\genfrac{}{}{0pt}{}{#1}{#2}\right)}
\newcommand{\ceil}[1]{\left\lceil #1 \right\rceil}
\newcommand{\floor}[1]{\left\lfloor #1 \right\rfloor}
\DeclareMathOperator{\lcm}{lcm}
\newcommand{\ZZ}{\mathbb{Z}}

\newcommand{\X}{\mathcal{X}}
\newcommand{\Y}{\mathcal{Y}}
\newcommand{\E}{\mathbb{E}}
\newcommand{\Variance}{\mathbb{VAR}}
\newcommand{\Real}{\mathbb{R}}
\newcommand{\expect}[1]{\mathbb{E}\left[#1\right]}
\newcommand{\proj}{\operatorname{proj}}
\newcommand{\Reg}{\mathrm{Reg}}
\newcommand{\inner}[2]{\left\langle #1,#2 \right\rangle}
\newcommand{\ind}[1]{\mathbf{1}_{\{#1\}}}
\newcommand{\norm}[1]{\left\lVert #1 \right\rVert}
\newcommand{\abs}[1]{\left\lvert #1 \right\rvert}

\newcommand{\Section}{Section}
\newcommand{\sref}[1]{\ref{#1}}
\newcommand{\Figure}{Figure}
\newcommand{\Figures}{Figures}
\newcommand{\fref}[1]{\ref{#1}}
\newcommand{\CaptionFontSize}{\small}
\newcommand{\BlankFootnote}[1]{%
  \begingroup
  \renewcommand{\thefootnote}{}%
  \footnote{#1}%
  \addtocounter{footnote}{-1}%
  \endgroup
}

% Give long mathematical prose a little extra flexibility at USC margins.
\setlength{\emergencystretch}{3em}

\theoremstyle{plain}
\newtheorem{Thm:Theorem}{Theorem}[chapter]
\newtheorem{Thm:Lemma}[Thm:Theorem]{Lemma}
\newtheorem{Thm:Assumption}[Thm:Theorem]{Assumption}
\newtheorem{Thm:Corollary}[Thm:Theorem]{Corollary}
\newtheorem{Thm:Remark}[Thm:Theorem]{Remark}
\newtheorem{Thm:Claim}[Thm:Theorem]{Claim}
\newtheorem{Thm:Proposition}[Thm:Theorem]{Proposition}
\newtheorem{Thm:Conjecture}[Thm:Theorem]{Conjecture}

\theoremstyle{definition}
\newtheorem{Thm:Problem}[Thm:Theorem]{Problem}
\newtheorem{Thm:Definition}[Thm:Theorem]{Definition}
\newtheorem{Thm:Example}[Thm:Theorem]{Example}
